\section{Validasi Form dengan JavaScript}

Validasi form di client memberikan umpan balik instan tanpa reload halaman \cite{mdn-learn-js}. Langkah umum: ambil nilai input dengan \texttt{value}, periksa kondisi (kosong, format email, panjang), tampilkan pesan error di elemen yang sesuai, dan panggil \texttt{event.preventDefault()} pada event \texttt{submit} jika validasi gagal \cite{mdn-learn-js}. Validasi client meningkatkan UX tetapi tidak menggantikan validasi server \cite{php-manual}. Gunakan Constraint Validation API (\texttt{checkValidity}, \texttt{setCustomValidity}) bila memungkinkan \cite{webdev}.
